% !TeX root = main.tex
% Writing guidance:
% - Describe the avalanche rescue workflow and alpine operational constraints.
% - Explain avalanche transceivers (ARVA/ARTVA), their capabilities and
%   limitations.
% - Distinguish established limitations from those the proposed approach might
%   address; do not assume it can replace existing rescue equipment.
\chapter{Avalanche Search and Rescue}

\begin{itemize}
\item SAR in mountainous regions introduction (alpine context)
\end{itemize}

\section{Avalanche Search and Rescue Practices}
\begin{itemize}
\item Practices and standard procedures in avalanche search and rescue operations
\item need to save the avalanche victims quickly to increase survival chances
\end{itemize}

\section{Avalanche Rescue Equipment and Technologies}
\begin{itemize}
\item Avalanche transceivers (ARVA/ARTVA), their capabilities and limitations
\item Bibliography on the use of artva in avalanche search and rescue operations
\item other tecnologies used in avalance SAR (drones, dogs, other sensors) and bibliography on this
\item Limitations of current technologies and potential areas for improvement including using different sensors such as LORA for coarse grained search
\item Potential integration of new sensor technologies with existing SAR workflows, which leads us to the next chapter: a background of LORA \& the general technical background for the thesis
\end{itemize}